%---------------------------------------------------------------------
% Title page and one orientation page.
%---------------------------------------------------------------------
\thispagestyle{empty}

\begin{tikzpicture}[remember picture, overlay]
  \fill[bmblued] (current page.north west) rectangle
        ([yshift=-12.3cm]current page.north east);
  \fill[bmaccent] ([yshift=-12.3cm]current page.north west) rectangle
        ([yshift=-12.5cm]current page.north east);
\end{tikzpicture}

\vspace*{0.9cm}
{\color{white}
  {\large علامہ اقبال اوپن یونیورسٹی، اسلام آباد}\par
  \vspace{3pt}
  {بی اے \ \textbullet\ \ ایسوسی ایٹ ڈگری}\par
  \vspace{20pt}
  {\Huge\bfseries سیرتِ طیبہ}\par
  \vspace{6pt}
  {\large\bfseries \textenglish{Seerat-e-Tayyaba}}\par
  \vspace{12pt}
  {\LARGE\bfseries متوقع سوالات اور ان کے جوابات}\par
  \vspace{12pt}
  {\large کورس کوڈ \textenglish{436}}\par
  \vspace{4pt}
  {سترہ سوالات \ \textbullet\ \ چار حقیقی پرچوں سے ماخوذ\par}
}

\vspace{1.0cm}

\begin{tcolorbox}[enhanced, colback=bmtint, colframe=bmblue, boxrule=0.6pt,
                  arc=3pt, left=12pt, right=12pt, top=10pt, bottom=10pt]
\textbf{\color{bmblue}سوالات کہاں سے لیے گئے ہیں؟}\ \
اس کورس کی مقررہ کتاب (\textenglish{787} صفحات) صرف تصویری اسکین ہے --- کوئی
متنی تہہ نہیں۔ اس لیے تمام سوالات \textbf{چار حقیقی گزشتہ پرچوں} سے لیے گئے
ہیں:

\smallskip
\begin{itemize}\setlength{\itemsep}{0pt}\setlength{\topsep}{2pt}
  \item بی اے --- سمسٹر \textbf{بہار \textenglish{2013}}، \textbf{خزاں
        \textenglish{2012}} اور \textbf{خزاں \textenglish{2017}}
  \item بی اے / ایسوسی ایٹ ڈگری --- سمسٹر \textbf{بہار \textenglish{2022}}
\end{itemize}

\smallskip
بتیس سوالات یکجا کر کے \textbf{سترہ موضوعات} بنے --- یہی دہراؤ بتاتا ہے کہ
کون سا موضوع سب سے اہم ہے۔

\medskip
\textbf{\color{bmblue}صحتِ مواد۔}\ \
تاریخیں، اعداد اور واقعات معیاری سیرت کی کتابوں سے تصدیق شدہ ہیں۔ جہاں کسی
پرچے کے مختصر سوال میں ایسی تفصیل پوچھی گئی جس کی تصدیق اس نشست میں نہ ہو
سکی، وہ سوال شامل نہیں کیا گیا --- اندازے سے نہیں لکھا گیا۔

\medskip
\textbf{\color{bmblue}حیثیت۔}\ \
ذاتی مطالعے کے لیے۔ یہ اے آئی او یو کی مطبوعہ کتاب نہیں۔ اسے من و عن نقل کر کے
اسائنمنٹ میں جمع کرانا منع ہے۔
\end{tcolorbox}

\clearpage

%---------------------------------------------------------------------
\section*{پرچے کا خاکہ}
\markboth{پرچے کا خاکہ}{}

پرچے کی ساخت چاروں پرچوں میں تقریباً ایک جیسی ہے، مگر ایک تبدیلی نظر انداز
نہیں ہونی چاہیے:

\medskip
\begin{center}
\begin{tabular}{@{}p{4.6cm} p{4.4cm} p{5.0cm}@{}}
\toprule
 & \textbf{\textenglish{2012--2017}}
 & \textbf{بہار \textenglish{2022}} \\
\midrule
کل نمبر            & \textenglish{100} & \textenglish{100} \\
کامیابی کے نمبر     & \textenglish{40}  & \textbf{\textenglish{50}} \\
وقت                & تین گھنٹے          & تین گھنٹے \\
سوالات دیے جاتے ہیں & آٹھ               & آٹھ \\
حل کرنے ہیں         & پانچ              & پانچ \\
فی سوال نمبر        & \textenglish{20}  & \textenglish{20} \\
\bottomrule
\end{tabular}
\end{center}

\medskip
\begin{ghalti}[سب سے مہنگی غلط فہمی]
کامیابی کے نمبر اب \textbf{چالیس نہیں، پچاس} ہیں --- پرانے پرچوں کی بنیاد پر
''چالیس کافی ہیں'' سمجھنا دس نمبر کے فرق سے فیل کرا دیتا ہے۔
\end{ghalti}

\medskip
\noindent
\textbf{\color{bmblue}کون سا موضوع کتنے پرچوں میں آیا}

\smallskip
\begin{center}
\begin{tabular}{@{}p{7.0cm} p{2.2cm} p{4.8cm}@{}}
\toprule
\textbf{موضوع} & \textbf{کتنے پرچوں میں} & \textbf{اس کتابچے میں} \\
\midrule
سیرت کے مآخذ و مصادر               & \textbf{چاروں} & سوال \textenglish{1} \\
نبی \saw کی کثیر الجہتی شخصیت (بحیثیت) & تین سے چار      & سوال \textenglish{13--16} \\
فتح مکہ                             & تین  & سوال \textenglish{11} \\
معاشی تعلیمات                       & تین  & سوال \textenglish{16} \\
غزوہ بدر اور میثاقِ مدینہ           & دو   & سوال \textenglish{7--8} \\
غزوہ خیبر                           & دو   & سوال \textenglish{10} \\
\bottomrule
\end{tabular}
\end{center}

\medskip
\begin{nukta}[سیرت کے جواب کا سانچہ]
ممتحن \emph{تاریخ، حوالہ اور ترتیبِ واقعات} دیکھتا ہے۔ ہر جواب میں:
\textbf{(۱)} مختصر تعارف، \textbf{(۲)} \textbf{تاریخ} (ہجری اور عیسوی)،
\textbf{(۳)} واقعات \emph{ترتیب سے}، \textbf{(۴)} نتائج و اسباق، اور جہاں
ممکن ہو \textbf{(۵)} متعلقہ آیت یا حدیث کا حوالہ۔ غزوات میں محض کہانی نہ
لکھیے --- \emph{اسباب، واقعات اور نتائج} تینوں الگ عنوانوں سے لکھنا اضافی
نمبر دلاتا ہے۔
\end{nukta}
